% --- MAGIC COMMENTS: CONFIGURAZIONE EDITOR --- (vedi main.tex per riferimenti e spiegazioni) ---
% !TEX root = ../../main.tex

% ____________________________________________ NEW CUSTOM MACROS (COMMANDS) ____________________________________________

% Di seguito vengono presentati i comandi custom/personalizzati dichiarati dall'utente.
% Vengono inoltre riepilogati i comandi custom definite nella classe di documento custom "HKNdocument".
% NOTA: Ricorda, ogni termine del tipo "\..." in LaTeX è una control sequence. Ogni control sequence è un Comando. I
% comandi possono essere primitivi o macro. Un primitivo (esempi \hline) è formulato direttamente nel sorgente del
% motore LaTeX. Una macro è invece una macroistruzione, cioè un insieme di altri comandi. L'utente in LaTeX dichiara
% solo macro, quindi ogni comando custom è precisamente una macro custom. Non può dichiarare "primitivi" (dovrebbe
% modificare il sorgente del motore grafico TeX). Ancora più precisamente oltre alle macro, l'utente può definire dei
% Registri (come i contatori \newcounter o le lunghezze \newlength), ma non sono comandi. Inoltre: su LaTeX non ci sono
% variabili, tuttavia si può usare una macro per contenere qualsiasi cosa, come fosse di fatto una variabile.

% NOTA BENE: COME SCRIVERE MACRO SU PIÙ RIGHE - ATTENZIONE AGLI A CAPO - TIPICI PROBLEMI DI FORMATTAZIONE
% Nella definizioe di macro, i simboli "%" prima degli a capo sono utili per evitare che vengano inseriti comandi di new
% line o di space (spazi bianchi in generale: un "a capo" diventa uno spazio, due o più "a capo" diventano un "a capo"
% nel pdf) indesiderati nel documento generato. Letteralmente vengono "commentati via". Come good pratice generale, è
% consigliato metterli sempre ogni qual volta si va a capo, in particolare quando si definiscono macro che generano
% blocchi di testo. Questo consenste di evitare in generale problemi di formattazione, particolarmente subdoli da
% identificare nel codice sorgente, e da risolvere). Altrimenti le Macro possono essere dichiarate tutte su una riga, ma
% è consigliabile solo per macro molto semplici e brevi.

% -------------------------------------- CAPITOLI INDICIZZATI MA NON NUMERATI ------------------------------------------

% COMANDO per creazione di CAPITOLI NON NUMERATI MA INDICIZZATI (APPAIONO NELL'INDICE), con corretto aggiornamento delle
% testatine (header, funzionante indipendentemente dalla programmazione della custom class HKNdocument, e senza problemi
% di compatibilità con essa).

% In generale NON USARE QUESTO COMANDO (consiglio valido per tutte le versioni presentate) ma preferisci \chapter
% nall'interno dell'attivazione di \frontmatter e \backmatter, qualora si debbano aggiungere capitoli non numerati
% ripsettivamente all'inizio e alla fine del documento. Per maggiori chiarimenti, leggi la PR 6 nel postambolo del
% main.tex, c'è una lunga spiegazione. Questo comando va usato solo per intermezzi del mainmatter, nel raro caso in cui
% non vadano numerati.

% NOTA: tutte le versioni presentati a seguire costituiscono dei wrapper, ossia del codice che avvolge \chapter* con i
% comandi correttivi necessari, per aggiungere le features mancanti (indicizzazione, aggiornamento header, link
% ipertestuale per l'indice in tutte le versioni, e il supporto agli short-title per le versioni 2 e 3, nativamente non
% supportato da chapter*)

% NOTA: in tutte le versioni, viene passato il titolo del chapter sia al leftmark (testatina pagine SX: di norma
% contiene il titolo del capitolo) che al rightmark (testatina pagine DX: di norma contiene il tittolo delle sezioni se
% ce ne sono, altrimenti rimane vuota), con markboth, per evitare le testatine di destra completamente vuote (un
% problema noto di LaTeX): infatti di norma non si usano le sezioni nei paragrafi non numerati (questa è la
% consuetudine), e il rightmark rimarrebbe vuoto (usando \markboth{#1}{} ad es.), lasciando così la testatina di destra
% vuota. Utilizzando \markboth invece il problema si risolve, e in modo flessibile: riempie la testatine di destra, ma
% qualora si aggiungessero sezioni, queste in automatico sovrascriverebbero il rightmark, lasciando così la testatina di
% destra con il giusto titolo di sezione, e non più con il titolo del capitolo. Tuttavia questo problema non si pone mai
% nella custom class HKNdocument (al momento della scrittura del commento: versione della class "2024/12/08 v1.5 Base
% class for publishing"), in quanto non stampa allo stato attuale il rightmark, ma solo il leftmark, in tutte le
% testatine. Sarebbe da correggere.

% VERSIONE 1: senza supporto agli short-title (dovrebbe essere il parametro opzionale)
% \newcommand{\uchapter}[1]{%
%     \chapter*{#1}%                          % Chiama capitolo non numerato non indicizzato
%     \phantomsection%                        % Aggiunge ancora per ipertesti (pacchetto hyperref), per indice
%     \addcontentsline{toc}{chapter}{#1}%     % aggiunge il capitolo all'indice
%     \markboth{#1}{#1}%                      % Aggiorna il contenuto della testatina
% }

% VERSIONE 2: supporto agli short-title in accordo al comando nativo \chapter[Short Title]{Long Title}, scritta
% sfruttando xparse, che offre direttive di alto livello per definire macro avanzate.
% Richiede per retrocompatibilità di aggiungere \RequirePackage{xparse} alla custom class, o \usepackage{xparse} nel
% preambolo del main.tex.

\NewDocumentCommand{\uchapter}{ o m }{%         % Dichiaratore di macro di xparse, più flessibile:
%                                               % o - optional: argomento opzionale (short title). Argomento #1.
%                                               % m - mandatory: argomento obbligatorio (long title). Argomento #2.
    \chapter*{#2}%                              % Chiamata a capitolo non numerato e non indicizzato
    \phantomsection%                            % Aggiunge ancora per ipertesti (pacchetto hyperref), per indice
    \addcontentsline{toc}{chapter}{#2}%         % Inserisce titolo LUNGO nell'indice
    \IfNoValueTF{#1}{%                          % Controllo Condizionale: controlla se è stato passato il titolo breve
        \markboth{#2}{#2}%                      % Se manca il titolo breve, usa il titolo lungo nella testatina
    }{%
        \markboth{#1}{#1}%                      % Se il titolo breve è presente, usa il titolo breve nella testatina
    }%
}

% VERSIONE 3: sfruttando solo direttive di basso livello (senza xparse)
% Vengono definiti in sostanza due comandi (motori) diversi, chiamati a seconda
% di come l'utente scrive il comando di chiamata \uchapter, che a livello di programmazione non è altro che un comando
% di facciata o dispatcher (similmente a quello che fa il comando \chapter, che chiama a sua vota \@chapter o \@schapter
% a seconda che sia usato con o senza asterisco ("star")).
% In realtà il 3° comando, ossia il Dispatcher, potrebbe essere messo come primo nella lista, essendo TeX un espansore
% di macro "pigro" (non controlla in anticipo se le macro inserite nelle definizioni esistano o meno: vengono compilate
% tutte insieme, poi se ci sono problemi, si verificano alla prima chiamata). Qui si è preferito rispettare l'ordine
% logico dei linguaggi più ferrei come il C, dove non si possono scrivere comandi prima di averli dichiarati.

% \makeatletter%                              % Rende la chiocchiala @, chiamata "at" in inglese, un carattere normale
% %                                           % rendendo possibile usarlo niella definizione di macro.
% %
% % 1. Sottomotore CON argomento opzionale (\uchapter[Breve]{Lungo})
% \def\@uchapter@withopt[#1]#2{%              % Definisce il comando da chiamare quando è presente l'arg. opzionale
% %                                           % #1 è il titolo breve (tra quadre, sia nella dichiarazione che nell'uso)
% %                                           % #2 è il titolo lungo (tra graffe nell'uso, senza graffe nella dichiaraz.)
%     \chapter*{#2}%                          % Chiamata a capitolo non numerato e non indicizzato
%     \phantomsection%                        % Aggiunge ancora per ipertesti (pacchetto hyperref), per indice
%     \addcontentsline{toc}{chapter}{#2}%     % Inserisce titolo LUNGO nell'indice
%     \markboth{#1}{#1}%                      % Usa il titolo BREVE nella testatina
% }

% % 2. Sottomotore SENZA argomento opzionale (\uchapter{Lungo})
% \def\@uchapter@noopt#1{%                    % Definisce il comando da chiamare quando è assente l'arg. opzionale
% %                                           % #1 è il titolo lungo (tra graffe nell'uso, senza graffe nella dichiaraz.)
%     \chapter*{#1}%                          % Chiamata a capitolo non numerato e non indicizzato
%     \phantomsection%                        % Aggiunge ancora per ipertesti (pacchetto hyperref), per indice
%     \addcontentsline{toc}{chapter}{#1}%     % Inserisce titolo LUNGO nell'indice
%     \markboth{#1}{#1}%                      % Usa il titolo LUNGO nella testatina
% }

% % 3. Il Dispatcher o "Semaforo"
% % \@ifnextchar guarda il carattere successivo nel testo dell'utente.
% \newcommand{\uchapter}{%                    % Definisce il comando da chiamare per generare capitoli indicizz. non num.
%     \@ifnextchar[%                          % Controllo condizionale: se al comando segue una parentesi quadra '['
%         {\@uchapter@withopt}%               % Chiama il comando con argomento opzionale: \@uchapter@withopt
%         {\@uchapter@noopt}%                 % Altrimenti, chiama il comando senza argomento opzionale: \@uchapter@noopt
% }

% \makeatother                                % Rende la chiocchiala @, chiamata "at" in inglese, di nuovo un carattere
% %                                           % speciale, rendendo impossibile usarlo nella definizione o chiamata di
% %                                           % macro (ripristina la condizione di default: @ è usata per proteggere
% %                                           % macro interne di basso livello, che non vanno ne chiamate ne modificate).


% ---------------------------- MACRO PER MONKEY PATCHING DEL BUG SUI CAPITOLO NON NUMERATI -----------------------------

% Per i dovuti chiarimenti, leggi il postambolo (nel main.tex) dopo ho parlato approfonditamente del bug dei capitoli
% non numerati, introdotti dalla custom class HKNdocument (Versione: 2024/12/08 v1.5 Base class for publishing).
% In particolare, il problema è affrontato nella sezione:
% "PR6: CAPITOLI NON NUMERATI - SOLUZIONI CON E SENZA MODIFICHE ALLA CUSTOM CLASS"
% Decommentando il codice a seguire, si ottiene del codice perfettamente funzionante per la versione sopra indicata
% della classe, che risolve il bug citato.

% \makeatletter
% \def\@chapter[#1]#2{%
%   \ifnum \c@secnumdepth >\m@ne
%     \if@mainmatter % <--- IL CONTROLLO MANCANTE REINSERITO
%       \refstepcounter{chapter}% Incrementa solo se siamo nel mainmatter
%       \typeout{\@chapapp\space\thechapter.}%
%       \addcontentsline{toc}{chapter}% Inietta nel ToC con il numero
%                        {\protect\numberline{\thechapter}#2}%
%     \else
%       \addcontentsline{toc}{chapter}{#2}% Inietta nel ToC SENZA numero
%     \fi
%   \else
%     \addcontentsline{toc}{chapter}{#2}%
%   \fi
%   \chaptermark{#1}% Aggiorna l'header con il titolo breve
%   \addtocontents{lof}{\protect\addvspace{10\p@}}%
%   \addtocontents{lot}{\protect\addvspace{10\p@}}%
%   \if@twocolumn
%     \@topnewpage[\@makechapterhead{#2}]%
%   \else
%     \@makechapterhead{#2}%
%     \@afterheading
%   \fi}
% \makeatother



% ---------------------------- CCUSTOM MACRO DEFINITI NELLA CUSTOM CLASS HKNdocument -----------------------------------

\begin{comment}     % --- RIEPILOGO MACRO (COMANDI) CUSTOM DEFINITI NELLA CLASSE "HKNdocument" ---

    RIEPILOGO MACRO (COMANDI) CUSTOM E AMBIENTI DEFINITI DALLA CLASSE HKNdocument E _preamble.tex
    NOTA: non possono essere ridichiarati con \newcommand{cmd}{def}, ma possono essere aggiornati con
    \renewcommand{cmd}{def}

    Il progetto HKN è stato programmato in modo da offrire alcune custom macro, grazie alla custom class HKNdocument
    (HKNdocument.cls, _preamble.tex, _postamble.tex). Alcune di esse funzionano in particolare anche grazie ai file di
    configurazione del progetto (latexmkrc, settings.json).

    --- INSERIMENTO IMMAGINI FACILITATO (da _preamble.tex) ---
    \addfigure[Didascalia]{nomefile}{0.8}       % Inserisce un'immagine posizionandola precisamente dov'è citata
                                                % (passa l'opzione [H]: exactly here) e centrandola, pescando il file
                                                % dalla cartella "res/". 0.8 sta per larga 80% della larghezzad del
                                                % testo (si può cambiare).
    \addsvg[Didascalia]{nomefile}{0.8}          % Inserisce un'immagine vettoriale SVG precisamente dov'è citata (passa
                                                % l'opzione [H]: exactly here) e centrandola, pescando il file dalla
                                                % cartella "res/svg/". 0.8 sta per larga 80% della larghezzad del testo
                                                % (si può cambiare).
    \adddrawio[Didascalia]{nomefile}{0.8}       % Inserisce uno schema grafico creato con Draw.io precisamente dov'è
                                                % citata (passa l'opzione [H]: exactly here) e centrandola, esportato in
                                                % SVG dalla cartella "res/drawio/". 0.8 sta per larga 80% della
                                                % larghezzad del testo (si può cambiare).

    NUOVA SPIEGAZIONE: Altro non sono che Wrapper dei comandi \includegraphics{} e \includesvg{} dei pacchetti graphicx
    e svg, pensati per facilitare l'inserimento delle immagini. Questi comandi inseriscono un'immagine posizionandola
    precisamente dove sono chiamati (passano l'opzione [H]: exactly here, del pacchetto float, più forte e robusta delle
    opzioni "h" o "h!" native), centrandola, pescando il file rispettivamente dalla cartella "res/"., "res/svg/" o
    "res/drawio/". (0.8 sta per larga 80% della larghezzad del testo, si può cambiare).
    Ecco, per questi comandi NON SERVE SPECIFICARE I PERCORSI! Li cercano in automatico nelle cartelle dedicate: basta
    specificare il nome del file, in particolare SENZA ESTENSIONE (le macro gestiscono automaticamente anche quella: per
    graphics la omettono, per svg l'aggiungono).

    --- BOX COLORATI PER LA TEORIA (da HKNdocument.cls) ---
    NOTA: questi riquadri speciali non sono generati a partire dagli ambienti offerti dal package amsthm, bensì sono una
    soluzione completamente personalizzata che si appoggia interamente al pacchetto grafico tcolorbox. Sono creati con
    la custom macro \DeclareHKNBox, citata a seguire.

    \begin{definition} ... \end{definition}     % Genera un riquadro con bordo ROSSO numerato per le Definizioni.
    \begin{theorem} ... \end{theorem}           % Genera un riquadro con bordo GIALLO numerato per i Teoremi.
    \begin{corollary} ... \end{corollary}       % Genera un riquadro con bordo ARANCIONE numerato per i Corollari.
    \begin{exercise} ... \end{exercise}         % Genera un riquadro con bordo BLU SCURO numerato per gli Esercizi.
    \begin{observation} ... \end{observation}   % Genera un riquadro con bordo GRIGIO-BLU numerato per le Osservazioni
                                                % generiche.

    --- MACRO "FABBRICA" PER CREARE NUOVI BOX COLORATI (da HKNdocument.cls) ---
    \DeclareHKNBox{nome}{TitITA}{TitENG}{colore}    % Macro avanzata. Usata per creare box colorati personalizzati
                                                    % (NOTA: nella nuova versione della classe sono stati aggiunti i
                                                    % lemmi).

    --- COMANDI TIPOGRAFICI SPECIALI (da HKNdocument.cls) ---
    \myuline{Testo}                             % Una sottolineatura intelligente che si interrompe elegantemente per
                                                % non tagliare le parti basse delle lettere (p, g, q, j). \myuline non è
                                                % mai utilizzato all'interno della custom class! Il bello è che per
                                                % inserire il pacchetto ulem che serve, si è creato il bug o conflitto
                                                % per la bibliografia. Vedi il main.tex per i rimandi.

    --- MACRO PER L'INSERIMENTO DEI METADATI (necessario usarle nel preambolo, per la custom class HKNdocument) ---
    \shorttitle{Titolo Breve}                   % Definisce un titolo accorciato da usare nelle intestazioni superiori
                                                % per non farle eventualmente sbordare.
    \editor{Nome Cognome}                       % Aggiunge la voce "Revisionato da: / Edited by:" nella pagina del
                                                % titolo.
    \docdate{Data Libera}                       % Imposta una data fissa o personalizzata per la copertina. Impostato
                                                % dal project manager scrivendo il semestre del corso.
    \docversion{1.0}                            % Aggiunge l'indicazione della versione del documento in copertina. 1.0
                                                % è un esempio di versione

    --- MACRO CUSTOM USATE COME VARIABILI PER OGGETTI DELLA PER COPERTINA E INTESTAZIONE DELLA CUSTOM CLASS HKNdocument
    (Modificabili con \renewcommand se necessario) (da HKNdocument.cls) ---
    \doclogofrontpage                           % Macro che contiene il logo principale della copertina
                                                % (Default: logo HKN).
    \doclogoupperpage                           % Macro che contiene il piccolo logo dell'intestazione
                                                % (Default: logo piccolo HKN).
    \doclicense                                 % Macro che contiene l'immagine della licenza
                                                % (Default: Creative Commons).
    \DDorganization                             % Macro che contiene il nome dell'organizzazione
                                                % (Default: IEEE-HKN Mu Nu Chapter).

    --- MACRO INTERNE DI SALVATAGGIO (BACKEND) USATE COME VARIABILI - ATTENZIONE: NON TOCCARE
    Queste macro rappresentano il "database interno" della custom class HKNdocument.

    \maintitle       % Salva il titolo esteso inserito dall'utente tramite il comando \title{...}.
    \DDauthors       % Raccoglie e concatena i nomi inseriti con le chiamate a \author{...}.
    \DDeditors       % Salva i nomi dei revisori inseriti tramite il comando custom \editor{...}.
    \DDdate          % Memorizza la data personalizzata definita tramite \docdate{...}.
    \DDversion       % Registra il numero di versione impostato con \docversion{...}.
    \STtitle         % Conserva il titolo breve definito con \shorttitle{...} per le intestazioni.

    REGOLE IMPORANTI (Perché non vanno modificate e perché non dovrebbero essere usate):
    1. NON SOCRASCRIVERLE (NO \renewcommand): Si romperebbe il collegamento logico con i comandi standard (FrontEnd)
    previsti dalla classe. Usa sempre i comandi ufficiali (\title, \author, ecc.) per inserire i dati.
    2. NON USARLE/RICHIAMARLE NEL TESTO: Queste "scatole" vengono lette dalla classe per costruire la copertina e le
    intestazioni. Spesso contengono regole di formattazione nascoste (es. font enormi e in grassetto). Stamparle in
    mezzo a un paragrafo distruggerebbe l'impaginazione. Inoltre, se una di queste macro fosse vuota, genererebbe un
    errore fatale di compilazione.


\end{comment}
